% Dirichlet Round 09 English TeX
% Translation of Werke II, pp. 121--138.
\documentclass[11pt]{article}
\usepackage[a4paper,margin=1in]{geometry}
\usepackage[T1]{fontenc}
\usepackage[utf8]{inputenc}
\usepackage{lmodern}
\usepackage{amsmath,amssymb}
\usepackage{microtype}
\setlength{\parindent}{1.2em}
\setlength{\parskip}{0.25em}
\newcommand{\J}[2]{\left(\frac{#1}{#2}\right)}
\begin{document}

\begin{center}
{\Large G. Lejeune Dirichlet's Works. II.}\par\vspace{1.2em}
{\Large XII.}\par\vspace{0.5em}
{\LARGE \textbf{On the First of the Proofs Given by Gauss\\ of the Reciprocity Law in the Theory of Quadratic Residues.}}\par\vspace{1em}
{\large By G. Lejeune Dirichlet.}\par\vspace{0.5em}
{\small Crelle, Journal für die reine und angewandte Mathematik, vol. 47, pp. 139--150.}
\end{center}

\begin{center}
{\large \textbf{On the First of the Proofs Given by Gauss\\ of the Reciprocity Law in the Theory of Quadratic Residues.}}
\end{center}

Among the numerous proofs of the fundamental theorem in the theory of quadratic congruences, the earliest, found by Gauss already in 1796 and published in the \emph{Disquisitiones arithmeticae} (Section IV),\footnote{Gauss' Works, vol. I, pp. 104--111. K.} has always seemed especially remarkable to me: both because of the very simple idea on which it rests, and because, so far as I know, this proof is the only one in which the argument never leaves the domain of congruences of the second degree, to which the theorem essentially belongs, whereas the other ways of grounding it rest on principles that appear more or less foreign to this domain.\footnote{Gauss himself judged his first proof in a later paper (Comment. soc. Gott. vol. XVI, p. 70) as follows: ``Sed omnes hae demonstrationes, etsi respectu rigoris nil desiderandum relinquent, quaeri videantur, e principiis nimis heterogeneis derivatae sunt, prima forsitan excepta, quae tamen per ratiocinia magis laboriosa procedit, operationibusque prolixioribus premitur.''} If this beautiful proof lacks the brevity by which some of the later ones are so highly distinguished, that defect is not inherent in the method. Rather it is due to the accidental circumstance that no notation suitable for calculation is used to express certain relations that recur frequently in this treatment; this made it necessary to distinguish eight different cases, each of which again splits into several subdivisions. By introducing the symbol first used by Legendre,\footnote{Gauss' Works, vol. II, p. 4. K.} in the more general meaning later given to it by Jacobi, and by making a few other simplifications which nevertheless do not alter the nature of the proof, the proof contracts so much that it seems scarcely inferior to any of the others in brevity, provided one does not ignore the fact that part of the developments occurring in it remains indispensable for the theory of quadratic residues even when another method of proof is chosen for the reciprocity law. I have thought myself justified in devoting a few pages to a new presentation of this subject simplified in the indicated way, all the more because experience has repeatedly shown me how much the understanding of the earlier proof is made difficult for beginning mathematicians by the large number of cases distinguished in it.

\section*{\S. 1.}

In this first paragraph a few elementary propositions and definitions needed below will be stated.

According as the congruence
\[
 x^2\equiv k\pmod m,
\]
in which $k$ denotes a positive or negative integer, is possible or impossible, $k$ is called a quadratic residue or non-residue of the modulus $m$, whose sign is of course immaterial. For our purpose it is permissible always to suppose that $k$ and $m$ have no common divisor. In the especially important case where $m$ is an odd prime $p$, the symbol
\[
 \J{k}{p}
\]
will denote the positive or negative unit according as $k$ is a quadratic residue or non-residue of $p$.

The well-known theorem that the product $k'k''\ldots$ is a quadratic residue or non-residue of $p$ according as those factors $k',k'',\ldots$ that are quadratic non-residues of $p$ occur in even or odd number is expressed, with the help of our notation, by
\[
 \J{k'k''\ldots}{p}=\J{k'}{p}\J{k''}{p}\cdots.
\]

For the possibility of the congruence
\[
 x^2\equiv k\pmod {p^{\varpi}},
\]
where $\varpi>1$, the condition
\[
 \J{k}{p}=1
\]
is plainly necessary; and one also easily proves that, once this condition is fulfilled, the congruence is soluble for every value of $\varpi$.

The congruence
\[
 x^2\equiv k\pmod {2^{\varpi}}
\]
behaves differently. For $\varpi=1$ the odd number $k$ has no condition to satisfy; if $\varpi=2$ or $\varpi\geq3$, the necessary and sufficient condition for solubility is respectively that $k$ have the form
\[
 4\varrho+1\quad\text{or}\quad 8\varrho+1.
\]

If the modulus in our congruence is the product of powers of different primes, then for its possibility it is necessary and sufficient that it be soluble modulo the several prime powers.

We shall now give the symbol $\J{k}{p}$, which was defined above for the case where $p$ denotes an odd prime not dividing the positive or negative number $k$, a more general meaning. Suppose that the odd number $m$, whose sign is immaterial, has no common divisor with $k$ and is decomposed into its equal or unequal prime factors $p',p'',p''',\ldots$, so that $m=p'p''p'''\cdots$. By the symbol $\J{k}{m}$ we shall understand the product
\[
 \J{k}{p'}\J{k}{p''}\J{k}{p'''}\cdots.
\]
It is easy to see that, in such a symbol, one may replace $k$ by any number congruent to it modulo $m$, and that for such symbols the two equations
\[
 \J{k}{m}\J{l}{m}=\J{kl}{m},\qquad
 \J{k}{m}\J{k}{m'}=\J{k}{mm'}
\]
hold.

But it must not be overlooked that $\J{k}{m}$ now no longer has the same relation to the congruence
\[
 x^2\equiv k\pmod m
\]
as before, when $m$ was prime. If the congruence is soluble, it still follows that $\J{k}{m}=1$, since then all the factors whose product defines $\J{k}{m}$ are equal to the positive unit. But the converse conclusion from the condition $\J{k}{m}=1$ to the possibility of the congruence is plainly not permitted.

Finally, still assuming $m$ odd, note that the possibility of the congruence
\[
 lx^2\equiv k\pmod m,
\]
when $k$ and $l$ are relatively prime to $m$, has as consequence the equation
\[
 \J{kl}{m}=1.
\]
This is immediate after bringing the congruence into the form
\[
 (lx)^2\equiv kl\pmod m.
\]

\section*{\S. 2.}

We now come to the proper subject of this paper: the criteria by which, for a given positive or negative number $k$, the simple odd moduli $p$ of which $k$ is a quadratic residue are distinguished from those with respect to which $k$ has the opposite behavior.

Since, by the theorem cited above, the required distinction reduces to the corresponding one for the prime factors of $k$, we have only to investigate the three cases in which $k$ is equal to one of the numbers $-1$, $2$, or to a positive odd prime $q$. For these three cases the desired criteria are contained in the following equations, where the odd prime $p$ is supposed different from $q$ and, like $q$, positive:
\[
\tag{a} \J{-1}{p}=(-1)^{\frac12(p-1)},
\]
\[
\tag{b} \J{2}{p}=(-1)^{\frac18(p^2-1)},
\]
\[
\tag{c} \J{q}{p}\J{p}{q}=(-1)^{\frac12(p-1)\frac12(q-1)}.
\]
According to the first of these equations, $\J{-1}{p}$ is $+1$ or $-1$ according as $p$ is contained in the form $4\mu+1$ or $4\mu+3$. According to the second, one has
\[
 \J{2}{p}=1\quad\text{for }p=8\mu+1,
 8\mu+7,
\]
and
\[
 \J{2}{p}=-1\quad\text{for }p=8\mu+3,
 8\mu+5.
\]
According to the third, one always has
\[
 \J{q}{p}=\J{p}{q},
\]
except when both primes $p,q$ have the form $4\mu+3$, in which case
\[
 \J{q}{p}=-\J{p}{q}.
\]

The equations $(a)$ and $(b)$ are easy to prove in certain special cases: for the first, in the case where $p$ has the form $4\mu+3$, so that $\J{-1}{p}$ must be $-1$. If this did not hold for all primes $4\mu+3$, let $p$ be the least for which $\J{-1}{p}=1$. Then one can set
\[
 e^2+1=ph,
\]
and if, as can always be done, $e<p$ is chosen even, then also $h<p$ and $h$ is of the form $4\mu+3$. Hence $h$ has a prime factor $r<p$ of the same form $4\mu+3$, for which the equation $e^2+1=ph$ gives
\[
 \J{-1}{r}=1,
\]
contrary to our supposition.

Secondly, to prove that always
\[
 \J{2}{p}=-1
\]
when $p$ is contained in one of the forms $8\mu+3$, $8\mu+5$, suppose that there are primes of one of these forms for which the theorem fails, and let $p$ be the least. Then one may again set
\[
 e^2-2=ph.
\]
If $e$ is chosen odd and also smaller than $p$, then, since
\[
 e^2-2=8\mu+7,
\]
the number $h$ will have the form $8\mu+5$ or $8\mu+3$ according as $p$ has the form $8\mu+3$ or $8\mu+5$, and it will also be smaller than $p$. Since no number of the form $8\mu+3$ or $8\mu+5$ can be formed from prime factors all contained in the forms $8\mu+1$, $8\mu+7$, there is therefore a prime factor $r$ of $h$ that has the form $8\mu+3$ or $8\mu+5$, and for which, by the preceding equation, contrary to our supposition, $\J{2}{r}=1$.

The proof, to be conducted in an entirely similar way, that
\[
 \J{-2}{p}=-1
\]
when $p$ has one of the forms $8\mu+5$, $8\mu+7$ will be omitted for brevity. If this latter result is put, for $p=8\mu+7$, into the form
\[
 \J{2}{p}\J{-1}{p}=-1,
\]
and if we note that, by what was proved above, since $8\mu+7$ is a special case of the form $4\mu+3$,
\[
 \J{-1}{p}=-1,
\]
then for primes $p=8\mu+7$ one obtains
\[
 \J{2}{p}=1.
\]
Thus theorem $(b)$ has now been proved for all primes not contained in the form $8\mu+1$.

\section*{\S. 3.}

Before we can undertake to prove generally the equations stated in the preceding paragraph, a few consequences are to be drawn from those equations, provisionally assumed correct. We first make the following observation.

If one wants to determine, for a product of odd factors $r$,
\[
 R=\Pi r,
\]
the remainder modulo $4$, and writes each factor $r$ in the form $(r-1)+1$, then in multiplying one may omit all terms in which first parts of these binomials are multiplied together. Thus
\[
 R\equiv1+\Sigma(r-1)\pmod4,
\]
or $\frac12(R-1)$ and $\Sigma\frac12(r-1)$ are of the same parity, that is, either both even or both odd.

In the same way, from $R^2=\Pi r^2$, since every odd square $r^2$ has the form $8\mu+1$, it follows that
\[
 R^2\equiv1+\Sigma(r^2-1)\pmod {64},
\]
and hence again that $\frac18(R^2-1)$ and $\Sigma\frac18(r^2-1)$ have the same parity.

With the help of this observation it is now easy to derive from the equations above the following more general equations of the same form, in which $P$ and $Q$ denote any two positive odd numbers with no common divisor:
\[
\tag{a'} \J{-1}{P}=(-1)^{\frac12(P-1)},
\]
\[
\tag{b'} \J{2}{P}=(-1)^{\frac18(P^2-1)},
\]
\[
\tag{c'} \J{Q}{P}\J{P}{Q}=(-1)^{\frac12(P-1)\frac12(Q-1)}.
\]
Putting $P=\Pi p$, where $p$ denotes each equal or unequal prime factor contained in $P$, applying theorem $(a)$ to each $p$, and multiplying all the resulting equations together, one obtains
\[
 \J{-1}{P}=(-1)^{\Sigma\frac12(p-1)},
\]
which passes into equation $(a')$ when the exponent is replaced by the number of the same parity $\frac12(P-1)$. The correctness of $(b')$ is seen in the same way.

To prove $(c')$, decompose also $Q$ into its simple factors $q$. The left-hand side can then be represented as a product of expressions of the form $\J{q}{p}\J{p}{q}$, where each $p$ is combined with each $q$. Replacing every such expression by its value given by formula $(c)$ gives
\[
 \J{Q}{P}\J{P}{Q}=(-1)^{\Sigma\frac12(p-1)\frac12(q-1)}.
\]
The summation sign here embraces all combinations $p,q$. The sum is therefore the product of the factors
\[
 \Sigma\frac12(p-1)\quad\text{and}\quad \Sigma\frac12(q-1),
\]
for which one may substitute the numbers of the same parity
\[
 \frac12(P-1)\quad\text{and}\quad \frac12(Q-1),
\]
and thus obtains $(c')$.

Multiplying $(a')$ and $(c')$ gives
\[
 \J{-Q}{P}\J{P}{Q}=(-1)^{\frac12(P-1)\frac12(Q+1)}
\]
or
\[
 \J{-Q}{P}\J{P}{-Q}=(-1)^{\frac12(P-1)\frac12(-Q-1)}.
\]
Thus equation $(c')$ remains valid when it is applied to the positive number $P$ and the negative number $-Q$, and in that case it is a consequence of $(a')$ and of the original equation $(c')$. Conversely, the latter is plainly a consequence of $(a')$ and of equation $(c')$ applied to the combination $P,-Q$. The same observation naturally applies also to equations $(a)$ and $(c)$, which are contained in those presently under consideration.

\section*{\S. 4.}

To prove the theorems $(a)$ and $(c)$ in general, suppose that both hold up to, but not including, an arbitrary positive odd prime $q$: that is, the first holds for every positive odd prime less than $q$, and the second for every two such primes. If one can derive from this assumption that $(a)$ holds for $q$, and that $(c)$ holds for every combination $p,q$ with $p$ a positive odd prime less than $q$, then both theorems, since they are plainly true for the first primes, are proved in general.

It is easy to see that the required verification reduces to the following two points.

First, if for a suitably chosen sign in $\varpi=\pm p$ one has
\[
 \J{\varpi}{q}=1,
\]
then, according to equation $(c)$ applied to the combination $q,\varpi$, one must have
\[
 \J{q}{\varpi}=\J{\varpi}{q}(-1)^{\frac12(q-1)\frac12(\varpi-1)}=(-1)^{\frac12(q-1)\frac12(\varpi-1)}.
\]

Second, when $q$ has the form $4\mu+1$ and $\J{p}{q}=-1$, it remains to derive $\J{q}{p}=-1$, and at the same time to show that this second case occurs for every $q$ of the form $4\mu+1$, that is, that among the primes $p<q$ there is always at least one which is a quadratic non-residue of $q$.

If $q$ has the form $4\mu+3$, then $(a)$ has already been proved for $q$; hence, by the observation at the end of the preceding paragraph, it is immaterial for which of the combinations $p,q$ and $-p,q$ one proves equation $(c)$. Since
\[
 \J{-1}{q}=-1,
\]
it follows that
\[
 \J{-p}{q}=-\J{p}{q},
\]
so one of these combinations is always contained in the first case.

If, on the other hand, $q$ has the form $4\mu+1$, then from the assumption
\[
 \J{p}{q}=1
\]
the first case (with $\varpi=p$) gives
\[
 \J{q}{p}=1,
\]
and from the assumption
\[
 \J{p}{q}=-1
\]
the second gives
\[
 \J{q}{p}=-1,
\]
as required. That moreover
\[
 \J{-1}{q}=1
\]
then follows as follows. For a prime $p$ belonging to the second case one has
\[
 \J{p}{q}=-1,
\qquad
 \J{q}{p}=-1.
\]
If now
\[
 \J{-1}{q}=-1,
\]
and consequently
\[
 \J{-p}{q}=1,
\]
then the first case would give
\[
 \J{q}{p}=1,
\]
which contradicts the result derived from the second case.

Before carrying out the two points just indicated, one must note that the hypothesis made at the beginning of this paragraph, together with the derivation in the preceding paragraph of $(a')$ and $(c')$ from $(a)$ and $(c)$, plainly implies the correctness of $(c')$ for any two odd relatively prime numbers, provided they are not both negative and all prime factors contained in them are less than $q$.

\section*{\S. 5.}

Under the condition of the first case,
\[
 \J{\varpi}{q}=1,
\]
one can set
\[
 e^2-\varpi=qf.
\]
In this equation, as can always be done, choose $e$ even and smaller than $q$; then $f$ is odd, positive\footnote{In what follows, Latin letters denote essentially positive numbers, whereas Greek letters denote numbers that may be positive or negative.} (since the numerical value of $\varpi$ is smaller than $q$), and also smaller than $q$. The equation requires a separate treatment according as $f$ is or is not divisible by $\varpi$.

I. If $f$ is not divisible by $\varpi$, then by \S. 1 one has
\[
 \J{\varpi}{f}=1,
\qquad
 \J{qf}{\varpi}=\J{q}{\varpi}\J{f}{\varpi}=1.
\]
Multiplication and application of $(c')$ to the combination $f,\varpi$ gives
\[
 \J{q}{\varpi}=(-1)^{\frac12(f-1)\frac12(\varpi-1)}.
\]
On the other hand, from the equation above, in which $e$ is even, it follows that
\[
 \frac12(f-1)\quad\text{and}\quad \frac12(q-1)+\frac12(\varpi+1)
\]
have the same parity. Substituting the latter for the former in the exponent and omitting $\frac14(\varpi^2-1)$ as even gives, as required,
\[
 \J{q}{\varpi}=(-1)^{\frac12(q-1)\frac12(\varpi-1)}.
\]

II. If $f$ contains the factor $\varpi$, put
\[
 e=\varpi\varepsilon,
\qquad f=\varpi\varphi,
\]
so that $\varpi$ and $\varphi$ have the same sign. From the resulting equation
\[
 \varpi\varepsilon^2-1=q\varphi
\]
one obtains
\[
 \J{\varpi}{\varphi}=1,
\qquad
 \J{-q\varphi}{\varpi}=\J{q}{\varpi}\J{-\varphi}{\varpi}=1.
\]
Multiplication and application of $(c')$ to the numbers $\varpi,-\varphi$, of which only one is negative, gives
\[
 \J{q}{\varpi}=(-1)^{\frac12(\varpi-1)\frac12(\varphi+1)},
\]
or, since $\frac12(q-1)$ and $\frac12(\varphi+1)$ plainly have the same parity,
\[
 \J{q}{\varpi}=(-1)^{\frac12(q-1)\frac12(\varpi-1)}.
\]

\section*{\S. 6.}

The hypothesis occurring in the second case,
\[
 \J{p}{q}=-1,
\qquad q=4\mu+1,
\]
does not, as in the first case, immediately allow one to write down an equation. One must first prove the existence of an auxiliary prime $p'$, smaller than $q$, satisfying
\[
 \J{q}{p'}=-1.
\]
If $q$ has the form $8\mu+5$, this presents no difficulty: then $q-2$ has the form $8\mu+3$ and consequently has a prime factor $p'$ of one of the forms $8\mu+3$, $8\mu+5$, smaller than $q$, for which $q\equiv2\pmod {p'}$ and hence $\J{q}{p'}=\J{2}{p'}$, or, by \S. 2, $\J{q}{p'}=-1$.

It is not so easy to show the existence of $p'$ when $q$ has the form $8\mu+1$. The need to provide this proof for that case was perhaps the greatest difficulty Gauss had to overcome in his first foundation of the reciprocity theorem. He reached it by a very sharp argument, which essentially amounts to the following.

Let $2m+1<q$, and suppose that $q$ is a quadratic residue of all odd primes not exceeding $2m+1$. Since $q\equiv8\mu+1$, by \S. 1 the congruence $x^2\equiv q$ is then soluble for every modulus which, besides an arbitrary power of $2$, contains only odd prime factors not exceeding $2m+1$. These conditions are satisfied by the product
\[
 1\cdot2\cdot3\cdots(2m+1)=M,
\]
and so one may set
\[
 k^2\equiv q\pmod M,
\]
with $k$ chosen positive. Then
\[
 (q-1^2)(q-2^2)\cdots(q-m^2)\equiv(k^2-1^2)(k^2-2^2)\cdots(k^2-m^2)\pmod M.
\]
Now the right-hand side, when multiplied by the factor $k$, which is relatively prime to $M$, can be written as the continuous product
\[
 (k+m)(k+m-1)\cdots(k-m).
\]
This product is a multiple of $M$, as is easily shown arithmetically, and also follows from the fact that, when divided by $M$, it represents a number of combinations. Hence the left-hand side must also be divisible by $M$. Writing the quotient of this division in the form
\[
 \frac1{m+1}\cdot \frac{q-1^2}{(m+1)^2-1^2}\cdot
 \frac{q-2^2}{(m+1)^2-2^2}\cdots
 \frac{q-m^2}{(m+1)^2-m^2},
\]
one obtains an obvious contradiction if $m$ is chosen to be the integer immediately below $\sqrt q$, since the quotient is then a product of proper fractions. It has tacitly been assumed that this choice of $m$ satisfies the condition $2m+1<q$ on which the argument is based. This is indeed the case, since $2m+1<2\sqrt q+1$, and $2\sqrt q+1$ is plainly smaller than $q$ for all primes $8\mu+1$, the least of which is $17$. Thus it is proved that there is always a prime $p'$ smaller than $2m+1$, and therefore smaller than $q$, for which
\[
 \J{q}{p'}=-1.
\]

Also note that for our auxiliary prime $p'$ one has
\[
 \J{p'}{q}=-1,
\]
because the assumption $\J{p'}{q}=1$ would, by the preceding paragraph, imply $\J{q}{p'}=1$.

We now pass to showing that from
\[
 \J{p}{q}=-1,
\]
where $q=4\mu+1$, it always follows that
\[
 \J{q}{p}=-1.
\]
We may regard $p$ as different from the auxiliary prime $p'$, since for the latter the simultaneous equations
\[
 \J{p'}{q}=-1,
\qquad
 \J{q}{p'}=-1
\]
have already been established. With their help, our hypothesis can be represented by
\[
 \J{pp'}{q}=1,
\]
and the conclusion to be drawn by
\[
 \J{q}{pp'}=1.
\]
From the first of these equations one can set
\[
 e^2-pp'=q\varphi,
\]
where $e$ is to be even and smaller than $q$. Then $\varphi$ is odd and, since $p<q$ and $p'<q$, numerically smaller than $q$. There are three cases according as $\varphi$ is divisible by none of the primes $p,p'$, by exactly one of them, or by both at once. Since the equation above and the relation to be derived from it are symmetric in $p$ and $p'$, it makes no difference in the treatment of the second case whether $p$ or $p'$ is taken as a factor of $\varphi$.

I. If $\varphi$ is divisible by neither $p$ nor $p'$, then from $e^2-pp'=q\varphi$ one obtains
\[
 \J{pp'}{\varphi}=1,
\qquad
 \J{q\varphi}{pp'}=\J{q}{pp'}\J{\varphi}{pp'}=1.
\]
Multiplication and application of $(c')$ give
\[
 \J{q}{pp'}=(-1)^{\frac12(pp'-1)\frac12(\varphi-1)}.
\]
But from the preceding equation, since $q=4\mu+1$, the numbers $\frac12(pp'-1)$ and $\frac12(\varphi-1)$ have opposite parity--one is even. Therefore
\[
 \J{q}{pp'}=1
\]
must hold.

II. By symmetry we may suppose $\varphi$ is divisible by $p'$. Put
\[
 \varphi=p'\psi,
\qquad e=p'g.
\]
Then $e^2-pp'=q\varphi$ becomes
\[
 p'g^2-p=q\psi,
\]
where $\psi$ is divisible by neither $p$ nor $p'$. From this last equation one gets
\[
 \J{pp'}{\psi}=1,
\qquad
 \J{p'q\psi}{p}=1,
\qquad
 \J{-p q\psi}{p'}=1.
\]
By multiplication,
\[
 \J{q}{pp'}=\J{pp'}{\psi}\J{\psi}{pp'}\J{p'}{p}\J{-p}{p'}.
\]
Applying $(c')$ to the two combinations $pp',\psi$ and $p',-p$ gives
\[
 \J{q}{pp'}=(-1)^{\frac12(pp'-1)\frac12(\psi-1)+\frac12(p+1)\frac12(p'-1)}.
\]
Replace $\frac12(\psi-1)$ by the number $\frac12(p+1)$, which has the same parity by the preceding equation, and replace $\frac12(pp'-1)$ by $\frac12(p-1)+\frac12(p'-1)$. Then the exponent has the value
\[
 \frac12(p+1)(p'-1)+\frac14(p^2-1),
\]
which is plainly even. Hence
\[
 \J{q}{pp'}=1.
\]

III. In the third case put
\[
 \varphi=pp'\psi,
\qquad e=pp'g.
\]
Then $e^2-pp'=q\varphi$ becomes
\[
 pp'g^2-1=q\psi,
\]
whence
\[
 \J{pp'}{\psi}=1,
\qquad
 \J{-q\psi}{pp'}=\J{q}{pp'}\J{-\psi}{pp'}=1,
\]
and then
\[
 \J{q}{pp'}=(-1)^{\frac12(pp'-1)\frac12(\psi+1)}.
\]
Since $\frac12(\psi+1)$ is plainly even, it follows again that
\[
 \J{q}{pp'}=1.
\]
Thus equations $(a)$ and $(c)$, and the equations $(a')$ and $(c')$ derived from them, have been proved in general.

\section*{\S. 7.}

It remains to fill the case of theorem $(b)$ left open above, for primes of the form $8\mu+1$.

Let $q$ denote any prime of this form, and suppose that the theorem already holds for all primes of the same form smaller than $q$, or, what by \S. 2 is the same thing, for all primes smaller than $q$. If from this hypothesis the equation $\J{2}{q}=1$ can be deduced, the theorem will hold without restriction. This equation will be proved by deriving a contradiction from the assumption $\J{2}{q}=-1$.

Choose an auxiliary prime $p$, smaller than $q$, such that
\[
 \J{p}{q}=-1.
\]
Then, under the assumption just made,
\[
 \J{2p}{q}=1,
\]
and consequently one may set
\[
 e^2-2p=q\varphi,
\]
where, taking $e$ odd and smaller than $q$, $\varphi$ will likewise be odd and, apart from sign, smaller than $q$. It remains to distinguish whether $\varphi$ is or is not divisible by $p$.

I. In the first case one immediately obtains
\[
 \J{2p}{q\varphi}=\J{2}{q}\J{2}{\varphi}\J{p}{q\varphi}=1,
\qquad
 \J{q\varphi}{p}=1,
\]
and then
\[
 \J{2}{q}=\J{2}{\varphi}\J{p}{q\varphi}\J{q\varphi}{p}
 =\J{2}{\varphi}(-1)^{\frac12(p-1)\frac12(q\varphi-1)}.
\]
Now equation $(b')$, in which the sign of $P$ is immaterial and which follows from $(b)$, is applicable to the number $\varphi$, all of whose prime factors satisfy equation $(b)$. Thus
\[
 \J{2}{\varphi}=(-1)^{\frac18(\varphi^2-1)}.
\]
Because of the assumption $\J{2}{q}=-1$, the last equation can therefore be written in the form
\[
 -1=(-1)^{\frac18[2(p-1)(q\varphi-1)+\varphi^2-1]}.
\]
The expression in brackets plainly changes by a multiple of $16$ if $q\varphi-1$ and $\varphi$ are replaced by other numbers congruent to them modulo $8$. But from $e^2-2p=q\varphi$, since
\[
 e^2\equiv1,
\qquad q\equiv1\pmod8,
\]
one immediately obtains
\[
 q\varphi-1\equiv-2p,
\qquad
 \varphi\equiv1-2p\pmod8.
\]
Thus the bracketed expression is congruent modulo $16$ to
\[
 -4p(p-1)+(1-2p)^2-1,
\]
hence to zero. Consequently the right-hand side of our equation is equal to the positive unit, in contradiction with the left-hand side.

II. If $\varphi$ is divisible by $p$, put
\[
 \varphi=p\psi,
\qquad e=pg,
\]
so that $e^2-2p=q\varphi$ becomes
\[
 pg^2-2=q\psi.
\]
From this equation and from the assumption $\J{2}{q}=-1$ one has
\[
 \J{2p}{q\psi}=-\J{2}{\psi}\J{p}{q\psi}=1,
\qquad
 \J{-2q\psi}{p}=\J{2}{p}\J{-q\psi}{p}=1.
\]
Then, as above,
\[
 -1=(-1)^{\frac18[2(p-1)(q\psi+1)+p^2+\psi^2-2]}.
\]
Replacing again $q\psi+1$ and $\psi$ by the numbers congruent to them modulo $8$ in consequence of the preceding equation, namely $p-1$ and $p-2$, one sees that the expression in brackets in the exponent is congruent modulo $16$ to
\[
 2(p-1)^2+p^2+(p-2)^2-2=4(p-1)^2,
\]
hence to zero. This gives the same contradiction as before.

\end{document}
